%\begin{frame}{Table of Contents}
%    \begin{columns}
%        \begin{column}{0.6\textwidth}
%            \begin{enumerate}
%                \item \textbf{Introduction and Goal}
%                \item \textbf{The Proof}
%                \begin{itemize}
%                    \item Khintchine's Theorem\\ (a.k.a. some mathematical curiosity)
%                    \item manipulating the max-stable property
%                    \item ``magical'' separation into three types
%                \end{itemize}
%                \item \textbf{Examples (FIXME: probably not, because too much?)}
%                \begin{itemize}
%                    \item uniform distribution
%                    \item gaussian distribution
%                \end{itemize}
%                \item \textbf{Summary}
%            \end{enumerate}
%        \end{column}
%        \begin{column}{0.4\textwidth}
%            \begin{center}
%                MAYBE\\SOME\\PICTURE
%            \end{center}
%        \end{column}
%    \end{columns}
%\end{frame}

%\begin{frame}{Extreme Value Theory -- Background}
%    \uncover<1->{Consider $N$ i.i.d. random variables $X_1,\dots,X_N$...}
%    \begin{columns}
%        \begin{column}{0.7\textwidth}
%            \begin{itemize}
%                \item[\bpt]<2-> Central Limit Theory $\longleftrightarrow$ $X_1 + \dots + X_N$
%                \item[\bpt]<3-> Extreme Value Theory $\longleftrightarrow$ $\max\klammer\left\{ X_1, \dots, X_N \right\}$
%            \end{itemize}
%        \end{column}
%        \begin{column}{0.3\textwidth}
%            \begin{center}
%                MAYBE\\SOME\\PICTURE
%            \end{center}
%        \end{column}
%    \end{columns}
%\end{frame}

%\begin{frame}{Extreme Value Theory -- Short Recap}
%    Consider $n$ i.i.d. random variables $X_1,\dots,X_n$
%    \begin{itemize}
%        \pause\item[\bpt] Central Limit Theory $\longleftrightarrow$ $Y := X_1 + \dots + X_n$
%        \begin{itemize}
%            \pause\item[] existence and shape of c.d.f. for $n\rightarrow\infty$?
%        \end{itemize}
%        \,\\
%        \pause\item[\bpt] Extreme Value Theory $\longleftrightarrow$ $M_n := \max\klammer\left\{ X_1, \dots, X_n \right\}$
%        \begin{itemize}
%            \pause\item[] existence and shape of c.d.f. for $n\rightarrow\infty$?
%            \pause\item[] Finite $n$
%            \item[] $P\kla\left( M_n \le x \right) = P\kla\left( X_1 \le x, \dots, X_n \le x \right) \pause = \left[ F(x) \right]^n$
%            \pause\item[] ... but $n\rightarrow \infty$ complicated
%        \end{itemize}
%    \end{itemize}
%\end{frame}

\begin{frame}{Extreme Value Theory -- Short Recap}
    Consider $n$ i.i.d. random variables $\xi_1, \dots, \xi_n$ with common c.d.f. $F(x)$\,\\ \,\\
    \begin{itemize}
        \pause\item[\bpt] Central Limit Theory: $\qquad \Sigma_n := \xi_1 + \dots + \xi_n$
        \begin{itemize}
            \pause\item[ ] $P(\Sigma_n \le x)$ for $n\rightarrow\infty$?
        \end{itemize}
        \,\\
        \pause\item[\bpt] Extreme Value Theory: $\qquad M_n := \max\klammer\left\{ \xi_1, \dots, \xi_n \right\}$
        \begin{itemize}
            \pause\item[ ] $P(M_n \le x)$ for $n\rightarrow\infty$?
        \end{itemize}
    \end{itemize}
    \,\\
%    \pause We can write down the c.d.f.
    {\pause}Finite $n$:
    \begin{itemize}
        \item[] $P\kla\left( M_n \le x \right) \pause = P\kla\left( \xi_1 \le x, \dots, \xi_n \le x \right) \pause = \left[ F(x) \right]^n$
    \end{itemize}
%    {\pause}... but $n\rightarrow \infty$ complicated
%    \begin{textblock}{0.35}(0.65, 0.25)
%        \begin{center}
%            MAYBE\\SOME\\PICTURE
%        \end{center}
%%        \begin{figure}[!ht]
%%            \centering
%%            \includegraphics[width=\linewidth]{../EE_Pictures/ev_distributions_simple.png}
%%        \end{figure}
%    \end{textblock}
\end{frame}

%\begin{frame}{Extreme Value Theory -- Background II}
%    \begin{tikzpicture}[remember picture,overlay]
%        \node[xshift=0cm,yshift=0cm] at (current page.north west) {
%		\includegraphics{../EE_Pictures/ev_distributions_simple.png}
%		};
%    \end{tikzpicture}
%\end{frame}

\begin{frame}{Why do we need scaling?}
%    Example: $F(x) = x$,\ \ $x\in [0, 1]$ (uniform distribution)\\
%    \,\\
%    {\pause}
    Without scaling
    \begin{itemize}
        \item[\bpt] $P(M_n \le x) = \left[ F(x) \right]^n \pause \longrightarrow 0$\quad ($F(x) < 1$)
    \end{itemize}
    \,\\ \,\\
    {\pause}Scaling:\quad $\tilde{M}_n := \frac{M_n - b_n}{a_n}$, ($a_n>0$)
    \begin{itemize}
        \pause\item[\bpt] $P(\tilde{M}_n \le x) = P\left( \frac{M_n - b_n}{a_n} \le x \right) \pause = P\big( M_n \le a_n x + b_n \big) \pause = \left[ F(a_nx + b_n) \right]^n \longrightarrow G(x)$
        \pause\item[\bpt] Example: For $F(x) = x$ choose $a_n = \frac{1}{n}$, $b_n = 1$
        \pause\item[\bpt] $P(\tilde{M}_n \le x) = \left[ F(a_nx + b_n) \right]^n = \left( 1 + \frac{x}{n} \right)^n \pause \longrightarrow \e^{x}$
    \end{itemize}
\end{frame}

%\begin{frame}{Why do we need scaling?}
%    Example: $F(x) = x$,\ \ $x\in [0, 1]$ (uniform distribution)\\
%    \,\\
%    {\pause}Without scaling
%    \begin{itemize}
%        \item[\bpt] $\left[ F(x) \right]^n = x^n \pause \longrightarrow 0$\quad ($x < 1$)
%%        \pause\item[\bpt] not very useful...
%    \end{itemize}
%    \,\\
%    {\pause}With scaling
%    \begin{itemize}
%        \item[\bpt] $\left[ F(a_nx + b_n) \right]^n$
%        \pause\item[\bpt] for $F(x) = x$ choose $a_n = \frac{1}{n}$ and $b_n = 1$
%        \pause\item[\bpt] $\left[ F(a_nx + b_n) \right]^n = \left( 1 + \frac{x}{n} \right)^n \pause \longrightarrow \e^{x}$
%    \end{itemize}
%\end{frame}